% Отдельное руководство: существующие схемы и скрипты не изменяются.
% Из корня проекта:
% .venv/bin/python scripts/render_phono_guide.py
% latexmk -xelatex -interaction=nonstopmode -halt-on-error -outdir=build/phono_guide docs/PHONO_GUIDE.tex
% Итоговый PDF можно скопировать из build/phono_guide/ в docs/.
\documentclass[a4paper,11pt]{article}
\PassOptionsToPackage{dvipsnames}{xcolor}
\input{preamble.tex}
\usepackage[margin=21mm]{geometry}
\usepackage{booktabs,tabularx,longtable,array,enumitem}
\hypersetup{pdftitle={Первый фонокорректор: от ВЧ-тракта к аудио},
  pdfsubject={Пояснения к AT-3600L, NE5532 и аналоговому питанию}}
\setlength{\parindent}{0pt}
\setlength{\parskip}{5pt plus 1pt}
\setlength{\emergencystretch}{2em}
\setlist{itemsep=2pt,topsep=4pt,leftmargin=*}
\renewcommand{\arraystretch}{1.16}
\newcolumntype{Y}{>{\raggedright\arraybackslash}X}
\newcolumntype{L}[1]{>{\raggedright\arraybackslash}p{#1}}
\setcounter{tocdepth}{2}
\newcommand{\Ohm}{\,\text{Ом}}
\newcommand{\kOhm}{\,\text{кОм}}
\newcommand{\uF}{\,\text{мкФ}}
\newcommand{\nF}{\,\text{нФ}}
\newcommand{\mkS}{\,\text{мкс}}
\newcommand{\millivolt}{\,\text{мВ}}
\newcommand{\file}[1]{\nolinkurl{#1}}
\newcommand{\practicalnote}[1]{\begin{quote}\small\textbf{Практический смысл. }#1\end{quote}}
\newcommand{\guidefigure}[2]{%
  \begin{figure}[htbp]\centering
  \IfFileExists{build/phono_guide/#1.pdf}{%
    \includegraphics[width=0.96\linewidth]{build/phono_guide/#1.pdf}%
  }{\fbox{\parbox{0.9\linewidth}{\small Для рисунка запустить из корня проекта:
    \file{.venv/bin/python scripts/render_phono_guide.py}.}}}
  \caption{#2}\end{figure}}
\ctikzset{amplifiers/fill=cyan!8,resistors/fill=white,grounds/scale=0.65,
  bipoles/length=1cm}

\begin{document}
\hypersetup{pageanchor=false}
\begin{titlepage}
\vspace*{17mm}
{\Huge\bfseries Первый фонокорректор:\\[3mm]от ВЧ-тракта к аудио\par}
\vspace{8mm}
{\Large AT-3600L, NE5532, балансный вход\\и питание на LM317/LM337\par}
\vspace{15mm}
{\large Подробное объяснение для инженера, который умеет проектировать
ВЧ-аналоговые тракты и впервые собирает аудиопредусилитель\par}
\vspace{9mm}
Это руководство объясняет уже существующие схемы проекта: зачем нужен каждый
узел, какие расчёты стоят за номиналами и что проверять на первом макете.
Основной вариант — три NE5532, балансный вход от головки и обычные линейные выходы.

Основные числовые примеры относятся к состоянию проекта на 25 ноября 2026 года:
182 Ом между плечами, два электролита по 1000 мкФ на канал, питание ±12 В.
Приведённые АЧХ рассчитаны; характеристики собранного устройства ещё не измерены.
Отдельный раздел дополнен результатами ngspice по снимку KiCad от 28 ноября 2026 года.
\vfill
\textbf{Как пользоваться.} Держать рядом схемы:\par
\file{ne5532_ad3600l_ballanced.tex};\par
\file{ne5532_at3600l.tex};\par
\file{analog_power_lm317_lm337.tex}.
В имени первого файла сохранено исходное написание; модель головки — AT-3600L.
\par\vspace{5mm}
Отдельный пояснительный документ. Существующие схемы, конфигурация и расчётные
скрипты остаются самостоятельными файлами.
\end{titlepage}
\hypersetup{pageanchor=true}

\tableofcontents
\newpage

\section{Что здесь знакомо ВЧ-инженеру, а что меняется}

Фонокорректор подчиняется той же теории линейных цепей, обратной связи и шума,
которую используют в ВЧ-трактах. Заново изучать электронику не требуется.
Меняются источник сигнала, диапазон частот и относительная важность погрешностей.
Напряжение головки измеряется милливольтами; полезный диапазон начинается с десятков
герц, а усиление на нижних частотах оказывается существенно выше, чем на 1 кГц.

\begin{tabularx}{\linewidth}{@{}L{0.22\linewidth}YY@{}}
\toprule
Привычный вопрос & В этом проекте & Практическое следствие\\
\midrule
Согласование порта & Головка работает на заданную высокоомную нагрузку &
47 кОм и ёмкость кабеля важны; терминатор 50 Ом здесь непригоден.\\
Полоса сигнала & Звуковая полоса, но ОУ остаётся мегагерцовым устройством &
АЧХ проверяется в аудиодиапазоне, генерация ищется и далеко выше него.\\
Спектр полезного сигнала & Базовая полоса и предыскажение при записи &
Нужно частотно-зависимое усиление RIAA.\\
Возвратный ток & На ВЧ важна индуктивность, на НЧ заметно и общее сопротивление &
Милливольт на общей земле уже сопоставим с сигналом головки.\\
Оценка шума & Импеданс головки комплексный и зависит от частоты &
Одного числа шум-фактора для источника 50 Ом недостаточно.\\
Генератор & Часто откалиброван на нагрузку 50 Ом &
Нужно измерять фактическую амплитуду на входе платы.\\
\bottomrule
\end{tabularx}

Фонокорректор не является усилителем мощности. Он не питает динамики и не должен
выдавать амперы. Его задача — подготовить слабый сигнал головки для входа AUX/LINE
или для последующего интерфейса АЦП. Поэтому здесь важнее шум, точность коррекции,
помехоустойчивость и запас по перегрузке, чем выходная мощность.

В руководстве три разных уровня утверждений. Назначение узлов и уравнения следуют
из схемы; числовые АЧХ получены идеальной и упрощённой моделями; шум, искажения,
реальный CMRR и устойчивость конкретной платы требуют измерений.
Эти уровни полезно держать раздельно: точный график расчёта не превращается
автоматически в паспорт устройства.

\section{Что представляет собой головка и почему здесь нет 50 Ом}

\subsection{Источник напряжения с сопротивлением и индуктивностью}

Для электрического анализа канал магнитной головки удобно представить
идеальным источником $E(s)$ с последовательным импедансом
\[
 Z_c(s)\approx R_c+sL_c.
\]
Это начальная модель. Реальная головка сложнее: присутствуют потери,
паразитные ёмкости и механическая система иглы. Значения $R_c$ и $L_c$ конкретной
AT-3600L в текущие скрипты не внесены. Приписывать ей произвольные параметры другой
MM-головки нельзя.

Если вход представляет собой $R_L\parallel C_L$, то напряжение на нём равно
\[
 V_{\mathrm{in}}(s)=E(s)\frac{Z_L(s)}{Z_c(s)+Z_L(s)},\qquad
 Z_L(s)=\left(\frac1{R_L}+sC_L\right)^{-1}.
\]
Изменение нагрузки может менять АЧХ даже идеального предусилителя. Индуктивность
головки и суммарная ёмкость образуют резонансную систему. Простейшая оценка
$f_0=1/(2\pi\sqrt{L_cC_L})$ показывает направление влияния, но не учитывает
демпфирование сопротивлениями.

\practicalnote{47 кОм — часть условий работы головки. Это не сопротивление для
максимальной передачи мощности. Короткое замыкание на 50 Ом изменит и уровень,
и электрическую нагрузку головки.}

\subsection{Что входит в нагрузочную ёмкость}

Суммарная ёмкость включает проводку тонарма, межблочный кабель, монтаж,
вход усилителя и отдельно установленный конденсатор. Площадка под конденсатор
на схеме обозначена DNP: по умолчанию деталь не устанавливается.
Подбор делают с учётом паспорта головки и уже имеющегося кабеля.

В балансном подключении нужно различать ёмкость между сигнальными проводами и
ёмкости каждого провода к экрану. Два одинаковых конденсатора к земле при
дифференциальном сигнале эквивалентны их последовательному соединению;
конденсатор непосредственно между входами действует всей своей ёмкостью.
Асимметрия ёмкостей к экрану дополнительно ухудшает баланс.

Головке AT-3600L не требуется питание от этой платы. На её выводы не подаётся
ни ±12 В, ни микрофонное фантомное питание. Цвета L+/L− и R+/R− обозначают
концы двух катушек, а не шины источника питания.

\section{Уровни сигнала, децибелы и коррекция RIAA}

\subsection{Милливольты превращаются в линейный уровень}

По конфигурации проекта расчётный пример — $4\millivolt$ действующего значения
на 1 кГц между выводами одной катушки. Это условие расчёта, а не полная
паспортная характеристика головки при любой записи. Номинальное напряжение
головок привязывается к частоте, скорости записи и условиям измерения.

Для коэффициента усиления по напряжению
\[
 G_{\mathrm{dB}}=20\log_{10}|G|.
\]
Усиление 100 В/В соответствует 40 дБ. Получаем примерно
$4\millivolt\cdot100=0.4$ В действующего значения.
Для синусоиды $V_{pp}=2\sqrt2\,V_{\mathrm{rms}}$:
0.4 В RMS — около 1.13 В от пика до пика. На осциллографе и в генераторе
нужно сверять, какое именно представление амплитуды выбрано.

В единицах dBV отсчёт ведётся от 1 В RMS; в dBu — примерно от 0.775 В RMS.
Эти логарифмические уровни не задают выходное сопротивление каскада.
Привычные dBm относятся к мощности и требуют известного сопротивления нагрузки.

\subsection{Почему усиление должно зависеть от частоты}

При записи пластинки вводят предыскажение: низкие частоты ослабляют относительно
средних, высокие усиливают. Воспроизведение требует обратной частотной коррекции.
Один широкополосный усилитель с коэффициентом 100 эту задачу не решает.

Стандартная передаточная функция RIAA-воспроизведения записывается как
\[
 H_{\mathrm{RIAA}}(s)=K\frac{1+s\tau_z}{(1+s\tau_1)(1+s\tau_2)},
\]
\[
 \tau_1=3180\mkS,\qquad \tau_z=318\mkS,\qquad \tau_2=75\mkS.
\]
Соответствующие частоты перегиба — примерно 50.05, 500.5 и 2122 Гц.
$K$ выбирают так, чтобы получить нужный уровень на 1 кГц.
При нормировке на 1 кГц эталон даёт примерно +19.27 дБ на 20 Гц,
+13.09 дБ на 100 Гц и −19.62 дБ на 20 кГц.

\guidefigure{response_comparison}{Расчётная АЧХ и отклонение от RIAA для двух
схем проекта. Верхний график нормирован на 1 кГц; нижний позволяет увидеть
малую погрешность, скрытую масштабом верхнего графика. Источник идеальный,
импеданс головки и кабель не включены.}

\subsection{Точность коррекции и субъективный результат}

Ошибка RIAA — это отклонение формы АЧХ от эталона после выбранной нормировки.
Если весь график поднялся на 0.2 дБ, это изменение усиления. Если поднялись
только высокие частоты, это изменение коррекции. Эти две погрешности желательно
измерять и показывать отдельно.

Не следует считать погрешность модели 0.01 дБ гарантией столь же точного
собранного корректора. На результат влияют допуски деталей, выходное сопротивление,
нагрузка, конечное усиление ОУ и сама головка. При этом согласование левого
и правого каналов полезно контролировать отдельно от общей ошибки относительно RIAA.

\subsection{RIAA и дополнительный инфранизкочастотный спад}

Сильное усиление низких частот захватывает рокот, коробление пластинки и
резонансы механической системы. Иногда вводят дополнительный НЧ-фильтр,
в том числе вариант RIAA/IEC с постоянной времени около 7950 мкс.
Это отдельное проектное решение, меняющее нижнюю часть АЧХ.

В локальном AND8177/D именно обсуждению такого спада посвящён раздел перед
рисунком 3\cite{onsemi-audio}. Поэтому отдельный конденсатор из этого рисунка
нельзя переносить в нашу схему с ожиданием прежней АЧХ. В наших расчётах эталон —
три постоянные времени RIAA; дополнительные разделительные цепи учтены явно.

\section{Что ОУ делает в этой схеме}

\subsection{От большого разомкнутого усиления к предсказуемому каскаду}

Для идеализированного ОУ
\[
 V_o=A(s)(V_+-V_-).
\]
В линейном режиме отрицательная обратная связь делает разность входных напряжений
малой. Отсюда привычное приближение $V_-\approx V_+$.
Оно работает, пока есть запас усиления, выход не ограничен и петля устойчива.
При насыщении или генерации подставлять равенство входов в расчёт уже нельзя.

У неинвертирующего каскада с $Z_f$ от выхода к инвертирующему входу
и $Z_g$ от этого входа к опорной точке
\[
 G(s)=1+\frac{Z_f(s)}{Z_g(s)}.
\]
Заменяя резисторы частотно-зависимыми импедансами, получаем активный фильтр.
Именно поэтому элементы обратной связи напрямую определяют слышимую АЧХ.

В NE5532 два усилителя, но один комплект питания. Выводы первой секции:
3 — неинвертирующий вход, 2 — инвертирующий, 1 — выход;
второй: 5, 6 и 7 соответственно. Вывод 8 — положительное питание,
вывод 4 — отрицательное\cite{ti-ne5532}.

\subsection{Смещение, входные токи и шум — разные параметры}

Даже при одинаковых внешних входных напряжениях реальный ОУ имеет эквивалентное
постоянное смещение $V_{OS}$. Входные токи $I_B$ создают дополнительные падения
на сопротивлениях. Например, условные 200 нА через 23.5 кОм дают 4.7 мВ.
Эта оценка объясняет порядок постоянных уровней; точный результат определяется
всеми ветвями и различием входных токов.

Шум — флуктуация во времени; смещение — постоянная или медленно меняющаяся ошибка.
Большой разделительный конденсатор не устраняет внутреннее смещение ОУ.
Он меняет усиление по постоянному току и таким образом ограничивает последствия.

Входному току нужен замкнутый путь. После разделительного конденсатора его
обеспечивает резистор к GND. Если убрать этот резистор, входной потенциал
не будет корректно определён\cite{adi-bias}. Катушка перед конденсатором
не заменяет такого пути по постоянному току.

\subsection{Аудиополоса не отменяет ВЧ-устойчивость}

Упрощённая модель в скриптах использует
\[
 A(s)=\frac{A_0}{1+sA_0/(2\pi\,GBW)},\qquad
 A_0=10^5,\quad GBW=12\,\text{МГц}.
\]
Она позволяет оценить потерю усиления и изменение АЧХ. Дополнительные полюса,
выходной импеданс и паразитные связи реальной микросхемы в неё не включены.
Поэтому этот расчёт не является проверкой запаса по фазе готовой платы.

Для ВЧ-инженера это знакомая ситуация: передаточная функция полезного сигнала
и устойчивость замкнутой системы — разные вопросы. На осциллографе нужно искать
высокочастотные колебания и при нулевом аудиосигнале.

\section{Исходный небалансный вариант: один ОУ на канал}

Файл \file{ne5532_at3600l.tex} использует один корпус NE5532 на два канала.
Сигнал L+ поступает через входной конденсатор, L− соединяется с местной входной
землёй. Для правого канала действует такая же схема.

На примере левого канала назначение деталей следующее:
\par
\begin{tabularx}{\linewidth}{@{}L{0.27\linewidth}Y@{}}
\toprule
Детали & Назначение\\
\midrule
$C_{A1}=2.2\uF$, $R_{A1}=47\kOhm$ & Разделение по постоянному току,
нагрузка головки и путь входному току ОУ.\\
$R_{A2}=90.9\kOhm$, $R_{A3}=10\kOhm$, $C_{A2}=31.5\nF$ &
Частотно-зависимая обратная связь, формирующая первую часть RIAA.\\
$R_{A4}=91\Ohm$, $C_{A3}=1000\uF$ BP &
Нижняя ветвь ООС: задаёт усиление для переменного сигнала и уменьшает
усиление постоянного смещения.\\
$R_{A5}=750\Ohm$, $C_{A4}=100\nF$ & Пассивная высокочастотная часть коррекции.\\
$C_{A5}=4.7\uF$, $R_{A6}=100\kOhm$ & Выходная развязка по постоянному
току и определённый потенциал выходного разъёма.\\
\bottomrule
\end{tabularx}

Импеданс верхней обратной связи
\[
 Z_f=R_f\parallel\left(R_z+\frac1{sC_f}\right),
 \quad R_f=90.9\kOhm,\ R_z=10\kOhm,\ C_f=31.5\nF.
\]
Он имеет два характерных масштаба времени, связанные с $R_zC_f$ и
$(R_f+R_z)C_f$. Полная АЧХ получается с учётом $1+Z_f/Z_g$,
а не только графика самого $Z_f$.

Пассивный выходной фильтр нагружается следующим устройством. При внешних 47 кОм
и собственных 100 кОм эквивалентная нагрузка составляет около 31.97 кОм.
Поэтому точный расчёт отличается от ненагруженного $1/(1+s\cdot75\mkS)$.
Текущий скрипт учитывает этот делитель и выходной разделительный конденсатор.

Получается около 99.13 В/В на 1 кГц и ошибка формы RIAA примерно
от −0.107 до +0.163 дБ в 20 Гц–20 кГц в идеальной модели.
Это простой и работоспособный принцип построения. Хорошая небалансная разводка
с короткой проводкой может дать вполне достойный результат.

\section{Балансный вариант: зачем три NE5532}

\subsection{Два конца катушки и отдельный экран}

В файле \file{ne5532_ad3600l_ballanced.tex} каждый канал имеет два входных
усилителя и один выходной разностный усилитель. Для стерео требуется шесть ОУ,
то есть три сдвоенных корпуса: U1 обслуживает левый вход, U2 — правый,
U3A и U3B формируют линейные выходы.

Катушка подключается между двумя сигнальными входами. Её напряжения удобно
разложить на дифференциальную и синфазную составляющие:
\[
 V_d=V_+-V_-,\qquad V_{cm}=\frac{V_++V_-}{2}.
\]
Сигнал, одинаковый на обоих проводах, относится к синфазному режиму.
Приёмник должен усиливать $V_d$ и подавлять $V_{cm}$.
Классическая трёхоперационная структура объясняется в MT-063\cite{adi-ina};
в нашем проекте её импедансы дополнительно задают коррекцию.

Входные резисторы — по 23.5 кОм к GND. Для дифференциального сигнала они
создают суммарную нагрузку 47 кОм между концами катушки. Каждый можно собрать
из двух параллельных 47 кОм. Два резистора по 47 кОм вместо них дали бы 94 кОм.
Входные 4.7 мкФ задают вместе с 23.5 кОм частоту около 1.44 Гц.

Экран и корпус тонарма подключаются отдельно. Нужно проверить существующую
проводку: L− и R− могут быть объединены или связаны с шасси в старом аппарате.
Проверяется и возможная связь контакта головки с её металлическими частями.
Одного разъёма с четырьмя контактами недостаточно для гарантии балансного включения.

\subsection{Откуда берётся коэффициент два}

Рассмотрим входную пару идеальных ОУ с одинаковыми $Z_f$ и ветвью $Z_g$
между инвертирующими входами. Отрицательная обратная связь удерживает
их потенциалы около соответствующих неинвертирующих входов.
Тогда ток через перемычку равен
\[
 I_g=\frac{V_+-V_-}{Z_g}=\frac{V_d}{Z_g}.
\]
Этот ток течёт через обе обратные связи. Получаем
\[
 X_+=V_++Z_fI_g,\qquad X_-=V_--Z_fI_g,
\]
\[
 X_+-X_-=\left(1+\frac{2Z_f}{Z_g}\right)V_d.
\]
Коэффициент два связан с вкладом двух выходов в их разность.
Он не означает, что напряжение катушки само удвоилось.
В симметричном возбуждении $+2$ и $-2$ мВ означают 4 мВ между концами катушки.
Подать по $+4$ и $-4$ мВ — уже получить 8 мВ дифференциального входа.

\subsection{Что происходит с общей помехой}

Если $V_+=V_-$, идеальный ток через $Z_g$ равен нулю.
Входные усилители повторяют общий сигнал с усилением около 1.
Разностный каскад затем его вычитает. Это позволяет входной паре иметь
большое дифференциальное усиление без столь же большого усиления синфазного сигнала.

Однако синфазное напряжение всё равно должно укладываться в допустимый входной
диапазон и размах выходов U1/U2. Если один из ОУ ограничился, последующее
вычитание не восстановит испорченный сигнал.

Балансное включение помогает от одинаково наведённых электрических помех и
ошибок общей опоры. Магнитное поле двигателя, наведённое непосредственно в
катушке, может создавать настоящий дифференциальный сигнал. Приёмник усилит его
вместе с музыкой. Поэтому баланс не заменяет размещение, экранирование и механику.

\subsection{Чего баланс сам по себе не обещает}

Собственные шумы двух входных ОУ складываются. Если учитывать только одинаковый
некоррелированный входной шум напряжения по 5 нВ/$\sqrt{\text{Гц}}$,
дифференциально получится около $\sqrt2\cdot5=7.1$ нВ/$\sqrt{\text{Гц}}$.
Это сравнение лишь одной составляющей, а не полного шума устройства.
Выигрыш балансного варианта заключается прежде всего в подавлении внешней
синфазной помехи. Итоговое соотношение сигнал/шум зависит от реальных условий.

\section{Зачем между плечами 182 Ом и два больших электролита}

\subsection{Две разные задачи одной ветви}

В нашем балансном канале
\[
 Z_g(s)=182\Ohm+\frac1{sC_a}+\frac1{sC_b},\qquad C_a=C_b=1000\uF.
\]
Малое сопротивление задаёт большое усиление для переменного сигнала и помогает
уменьшить вклад шума этой части схемы. Конденсаторы делают импеданс ветви
большим по постоянному току. В результате при $s\to0$ дифференциальное усиление
входной пары стремится к 1, а не к $1+2R_f/R_g\approx1000$.

Это важно даже при входных разделительных конденсаторах. Они блокируют внешнее
постоянное напряжение, но не устраняют $V_{OS}$ и падения от входных токов
самих ОУ. Если убрать конденсаторы из межплечевой ветви, разность этих ошибок
может сильно усилиться и отнять запас выходного размаха.

\practicalnote{Большие электролиты находятся в цепи, задающей усиление. Их назначение —
развести режимы DC и AC. Это не резервуар питания и не требование головки
иметь именно такую ёмкость.}

\subsection{Как получается неполярная составная ёмкость}

Два обычных алюминиевых электролита включены встречно: их отрицательные выводы
соединены друг с другом, а положительные смотрят к двум внешним узлам.
\begin{figure}[htbp]\centering
\begin{circuitikz}[font=\small]
  \draw (0,0) node[ocirc] {} node[above left] {к U1A(−)}
    to[R,l=$182\Ohm$] (2,0)
    to[C,l=$C_a\;1000\uF$] (5,0) node[circ] {}
    to[C,l=$C_b\;1000\uF$] (8,0) -- (9.2,0)
    node[ocirc] {} node[above right] {к U1B(−)};
  \node at (3.15,-0.4) {$+$};
  \node at (6.85,-0.4) {$+$};
  \node[align=center] at (5,-1.05) {соединённые минусы;\\плавающая точка};
\end{circuitikz}
\caption{Упрощённое изображение межплечевой ветви. Средняя точка конденсаторов
не соединяется ни с GND, ни с отрицательной шиной питания.}
\end{figure}

В линейном расчёте
\[
 C_{\mathrm{экв}}=\frac{C_aC_b}{C_a+C_b}=500\uF.
\]
Такое соединение описано в руководстве Cornell Dubilier\cite{cde}.
Это практический способ получить требуемую функцию
без поиска отдельного крупного BP-конденсатора.
В проекте выбраны алюминиевые электролиты с жидким электролитом по 1000 мкФ / 25 В.
На электролиты с проводящим полимером и танталовые детали такой выбор
автоматически не переносится.

В малосигнальной AC-модели полярность не участвует в уравнениях: там два идеальных
конденсатора. Утечки, электрохимическая асимметрия, ESR и нелинейность реальных
деталей таким расчётом не проверяются. Схематические знаки «+» нужны монтажу,
а не линейному решателю.

25 В здесь являются рабочим номиналом каждой отдельной детали, а не обещанием
выдержать 50 В на всей паре. В нормальном режиме напряжение на ветви мало:
оно определяется сигналом и разностью смещений входов. При включении возможны
дополнительные переходные процессы, поэтому подключение усилителя мощности
и громкости откладывают до первичной проверки платы.

\subsection{Почему 100 мкФ — уже другой результат}

Для большого дифференциального усиления характерный масштаб дополнительного
низкочастотного спада оценивается как
\[
 f_g\approx\frac1{2\pi R_gC_{\mathrm{экв}}}.
\]
Это частота отдельной ветви, а не полное описание всей многозвенной АЧХ.
Точный результат получаем из полной передаточной функции.

\begin{center}
\begin{tabular}{@{}rrr@{}}
\toprule
$C_{\mathrm{экв}}$, мкФ & $f_g$, Гц & Ошибка RIAA на 20 Гц, дБ\\
\midrule
500 & 1.75 & −0.068\\
100 & 8.75 & −0.794\\
50 & 17.49 & −2.499\\
\bottomrule
\end{tabular}
\end{center}
Ошибки указаны для нашей полной идеальной схемы с нормировкой на 1 кГц.
Один неполярный 100 мкФ даёт 100 мкФ. Два по 100 мкФ последовательно дадут
50 мкФ. Для неполярного керамического конденсатора встречно-последовательная
пара сама по себе не требуется.

На 20 Гц модуль сопротивления 500 мкФ — около 15.9 Ом, на 1 кГц — 0.318 Ом.
На фоне 182 Ом это даёт небольшой дополнительный сдвиг. Если условно поддерживать
4 мВ RMS между входами на 20 Гц, на всей идеальной ёмкости будет около 0.35 мВ RMS.
Это объясняет, почему применение здесь конденсаторов на 25 В не означает,
что в нормальной работе они находятся под напряжением питания.

\guidefigure{bridge_capacitance}{Замена только межплечевой ёмкости при неизменных
182 Ом и прочих элементах. Малые ёмкости увеличивают спад на нижнем краю.
Это расчёт АЧХ, а не измерение нелинейности конденсаторов.}

\subsection{Можно ли получить меньшую ёмкость корректным пересчётом}

Можно масштабировать импедансы: увеличить сопротивления входной ООС в $k$ раз,
а её ёмкости уменьшить в $k$ раз. Отношения импедансов и RC-произведения
при этом сохранятся. Для точного масштабирования $k=5$ получаются
910 Ом вместо 182 Ом, 454.5 кОм вместо 90.9 кОм, 50 кОм вместо 10 кОм,
6.3 нФ вместо 31.5 нФ и 100 мкФ вместо 500 мкФ.

Это допустимый инженерный вариант, но он повышает сопротивления шумящей сети.
Практический близкий набор 910 Ом / 453 кОм / 49.9 кОм / 6.3 нФ и два
электролита по 220 мкФ тоже был рассчитан: около 100.55 В/В на 1 кГц,
ошибка RIAA примерно от −0.073 до +0.012 дБ.
Эти числа иллюстрируют компромисс. В основной схеме сохранены 182 Ом и
два по 1000 мкФ, поскольку ограничение на их размер снято.

\section{Как U3 одновременно вычитает сигналы и завершает RIAA}

\subsection{Обычный разностный усилитель с частотно-зависимой обратной связью}

У каждого выходного каскада четыре резистора по 10 кОм и два одинаковых
конденсатора по 7.5 нФ. Положительное плечо приходит на неинвертирующий вход
через резистор; отрицательное — на инвертирующий. Две согласованные ветви
имеют импеданс
\[
 Z(s)=R\parallel\frac1{sC}=\frac{R}{1+sRC}.
\]
Делитель положительного входа задаёт
$V_p=X_+Z/(R+Z)$. Для идеального ОУ выход равен
\[
 V_o=\left(1+\frac ZR\right)V_p-\frac ZR X_-
     =\frac ZR(X_+-X_-).
\]
Следовательно,
\[
 H_2(s)=\frac1{1+sRC},\qquad RC=10\kOhm\cdot7.5\nF=75\mkS.
\]
Один ОУ выполняет и преобразование дифференциального сигнала в обычный,
и высокочастотную часть коррекции. Это сохраняет число корпусов равным трём.

Именно поэтому на выход балансной схемы не переносится прежний фильтр
750 Ом / 100 нФ. Его повторная установка добавила бы ещё один полюс 75 мкс
и изменила высокочастотный спад.

\subsection{Какие элементы должны совпадать}

Условия вычитания относятся к отношениям сопротивлений и импедансов.
Если один резистор 10 кОм отличается от остальных, часть синфазного сигнала
попадёт на выход. Если отличаются конденсаторы пары 7.5 нФ, ошибка будет
зависеть от частоты. Поэтому полезны согласованные резисторные сборки
и подбор конденсаторов внутри пары.

Абсолютный допуск и согласование — разные характеристики. Две ёмкости
с одинаковым отклонением +1\% хорошо согласованы, но обе сдвигают полюс RIAA.
Стабильный температурный дрейф отношений также важен. Надпись 0.1\% на
отдельных резисторах сама по себе не задаёт измеренный CMRR всего устройства.

\subsection{Зачем после U3 ещё резистор, конденсатор и резистор}

100 Ом после выхода отделяют ОУ от ёмкостной нагрузки кабеля и ограничивают
некоторые переходные токи. Обратная связь взята до этого резистора:
кабель не входит в неё напрямую. Полезная нагрузка от 47 кОм делает потерю
уровня на 100 Ом очень небольшой.

4.7 мкФ блокируют постоянную составляющую перед выходным разъёмом.
100 кОм после них задают нулевой потенциал без подключённого следующего
устройства и обеспечивают путь разряда. При внешних 47 кОм ёмкость видит
эффективную нагрузку около 31.97 кОм; низкочастотный масштаб этой выходной
цепи близок к 1.1 Гц.

Входы AUX/LINE рассчитаны на напряжение относительно земли. Вход АЦП может
иметь другое сопротивление, допустимый размах и постоянное смещение.
Его интерфейс проверяют отдельно: наличие разделительного конденсатора
не делает выход совместимым с любым АЦП.

\clearpage
\section{Сохранённый AC-анализ ngspice от 28 ноября 2026 года}

Это проверка передаточной функции по соединениям и номиналам из сохранённого
экспорта KiCad. Опубликованные данные не требуют наличия исходного каталога
\file{build/ngspice/run-_ttt3q4u}: снимок сохранён в
\file{docs/simulation/ngspice-2026-11-28/} вместе с CSV, комплексными отсчётами
DAT, SPICE-нетлистами и конфигурацией.

\input{ngspice_results.tex}

\begin{figure}[htbp]
\centering
\includegraphics[width=0.96\linewidth]{ngspice-2026-11-28-response-ru.pdf}
\caption{ngspice: усиление, фаза и нормированная ошибка RIAA.
Сплошная линия L и пунктирная R практически совпадают для номинально
одинаковых каналов. Вертикальная линия отмечает опорную частоту 1 кГц.}
\end{figure}

\clearpage
\section{Запас по уровню: проверяется каждый каскад}

Число «40 дБ» относится к усилению полного канала на 1 кГц.
На 20 Гц оно близко к 59.27 дБ, то есть примерно 919 В/В.
Одинаковые 4 мВ тестового синуса на всех частотах дали бы около 3.68 В RMS
на 20 Гц. Поэтому контроль АЧХ начинают с существенно меньшего уровня.

Кроме того, U1/U2 и U3 имеют разное частотное усиление.
На 1 кГц коэффициент между дифференциальным входом и разностью внутренних
выходов около 111.7, а полный коэффициент около 100.7.
На 20 кГц внутренняя разность всё ещё усиливается примерно в 100 раз,
а полный канал — примерно в 10.52 раза.

\guidefigure{stage_gains}{Распределение усиления в балансном канале. Кривая
входной пары относится к разности $X_+-X_-$, а не к каждому выходу по отдельности.
Затухание U3 на высоких частотах не уменьшает уже достигнутый размах U1/U2.}

При симметричном входе каждый из первых выходов содержит примерно половину
дифференциального выходного сигнала плюс общую составляющую.
Нужно контролировать оба относительно GND. Малое напряжение после U3 ещё не
доказывает отсутствие перегрузки раньше него: высокочастотный фильтр может
ослабить уже искажённый сигнал.

Допустимый размах не равен напряжению шин питания. Он зависит от нагрузки,
частоты и конкретного ОУ. При ±12 В нельзя автоматически ожидать линейный
выход от −12 до +12 В. Границу находят по паспорту в подходящих условиях
и затем по измерению своего макета.

Случайные щелчки пластинки, инфранизкочастотные колебания и сильная общая помеха
могут проверять разные ограничения схемы. Для измерения запаса повышают
входной уровень на выбранной частоте и смотрят все каскады. В дБ запас равен
$20\log_{10}(V_{\mathrm{in,max}}/V_{\mathrm{in,ref}})$ при одинаковом определении
амплитуды. Границу фиксируют по оговорённому критерию искажений либо видимого
ограничения, а не только по появлению плоской вершины.

\section{Шумы: что нужно считать и что можно измерить}

\subsection{Три основных вида внутренних источников}

Для анализа нужны как минимум входной шум напряжения ОУ, его входные шумовые
токи и тепловой шум сопротивлений. Они передаются к выходу через разные
частотно-зависимые функции\cite{adi-noise}.

Шумовое напряжение резистора при температуре $T$ имеет спектральную плотность
\[
 e_R=\sqrt{4kTR}.
\]
При $T=300$ К для 182 Ом получаем около 1.74 нВ/$\sqrt{\text{Гц}}$,
для 910 Ом — 3.88 нВ/$\sqrt{\text{Гц}}$.
Это шум самого резистора, а не полный шум предусилителя. Точный вклад на выходе
включает передаточную функцию его позиции в сети.

Шумовой ток даёт напряжение на импедансе источника. Например, условные
0.7 пА/$\sqrt{\text{Гц}}$ на 10 кОм соответствуют 7 нВ/$\sqrt{\text{Гц}}$.
Числа служат примером масштаба. У головки $Z_c(f)$ меняется с частотой,
поэтому один резистор в качестве её эквивалента не воспроизводит все условия.

Входные 47 кОм тоже шумят, но их полное ненагруженное шумовое напряжение
не прикладывается к входу целиком: его изменяет импеданс головки.
Нельзя просто сложить паспортные 5 нВ ОУ и шум разомкнутого резистора,
не учитывая сеть.

\subsection{Почему 5 нВ на корень герца не дают готовый SNR}

Чтобы получить действующее шумовое напряжение, нужно интегрировать
спектральную плотность мощности по полосе с учётом передаточных функций:
\[
 v_{n,o}^2=\int_{f_1}^{f_2}S_{n,\mathrm{in}}(f)|H(f)|^2\,df.
\]
Эта запись применима, когда источники уже приведены ко входу с корректными
частотными весами. Для самостоятельных источников внутри схемы удобнее
суммировать их вклады на выходе, каждый со своей функцией передачи.

Низкочастотное усиление RIAA велико, а у ОУ существует область шума $1/f$.
Поэтому оценка через усиление на 1 кГц и ширину 20 кГц может оказаться неверной.
Измеренный шум также зависит от выбранной полосы, фильтра взвешивания,
оконной функции и способа усреднения FFT.

Сравнивая два результата, записывают: подключённый источник или эквивалент,
диапазон частот, наличие A-взвешивания, усиление на 1 кГц и единицы RMS.
Шум с закороченным входом характеризует важный частный случай, но не полностью
описывает работу с индуктивной MM-головкой.

\subsection{Внешняя помеха — отдельная задача}

Для диагностики полезно различать шипение, сетевые линии и высокочастотную
генерацию. Линия 50 Гц может указывать на магнитную или электрическую наводку;
100 Гц часто связаны с выпрямленным питанием. Это подсказки, а не однозначный
диагноз. Гармоники могут появляться от нескольких разных механизмов.

ВЧ-сигнал способен детектироваться входными нелинейностями и проявляться в
аудиополосе. Для инженера импульсных трактов это знакомый механизм преобразования
помехи. Шумный цифровой преобразователь, USB-соединение и петли земли следует
проверять как реальные пути связи, а не оценивать только по названию стабилизатора.

\section{Компоненты: где нужна точность, а где большой корпус допустим}

\begin{longtable}{@{}L{0.30\linewidth}L{0.64\linewidth}@{}}
\toprule
Позиция & Критерий выбора\\
\midrule\endhead
31.5 нФ в ООС входной пары & Плёнка или подходящий C0G. Важны абсолютная
ёмкость, согласование плеч и температурная стабильность.\\
7.5 нФ в U3 & Плёнка или C0G; особенно важно согласование двух RC-ветвей
внутри канала.\\
Входные/выходные 4.7 мкФ & В проекте плёночные. При замене другим типом
учитываются утечки, нелинейность и фактическая ёмкость.\\
Два по 1000 мкФ между плечами & Обычные алюминиевые электролиты, одинаковая
серия и номинал удобны. Знаки полярности относятся к каждой детали отдельно.\\
100 нФ по питанию & Керамика X7R подходит для локальной развязки;
короткое соединение с выводами и землёй важнее большого расстояния до общего фильтра.\\
47 мкФ локально по питанию & Алюминиевые электролиты. На отрицательной шине
положительный вывод обращён к GND.\\
Резисторы сигнальной части & Металлоплёночные/тонкоплёночные. Точность отношений
и низкий избыточный шум важнее большой паспортной мощности.\\
Резисторы CRC-фильтра & Нужны заданное сопротивление, запас по рассеиванию
и подходящее исполнение для питания. Это другая задача, чем подбор RIAA-сети.\\
\bottomrule
\end{longtable}

Плёночные и C0G-конденсаторы удобны там, где нужна линейность частотозадающей
цепи. Измерения TI показывают различия между типами в сигнальном применении,
но их конкретные уровни искажений нельзя переносить на нашу другую схему
и напряжения без проверки\cite{ti-caps}.

Большая керамика X5R/X7R компактна, однако её ёмкость зависит от режима;
ферроэлектрическая керамика также обладает электромеханическим эффектом\cite{murata}.
Поэтому 100 мкФ в корпусе 2220 нужно оценивать по конкретной маркировке
и условиям работы. На межплечевой ветви постоянное напряжение мало;
потерю ёмкости, измеренную для питания под несколькими вольтами DC,
нельзя автоматически приписывать этому месту.

Если две детали имеют одинаковую надпись, это ещё не согласованная пара.
При измерении LCR-метром следует использовать одинаковые частоту, амплитуду
и схему замещения. Это особенно существенно для больших керамических ёмкостей.
У электролитов ±20\% допуска и заметный ESR — нормальные возможные паспортные
условия; точность НЧ-полюса оценивают с учётом этого допуска.

В исходной небалансной схеме обозначен отдельный 1000 мкФ BP. Если вместо него
применять пару по 2200 мкФ, линейная эквивалентная ёмкость будет 1100 мкФ,
а не 1000 мкФ. Это пример расчёта альтернативы; исходный файл и его конфигурация
такой заменой автоматически не обновляются.

Замена NE5532 более быстрым или менее шумным ОУ требует проверки всей схемы.
Распиновка, минимальное устойчивое усиление, входные токи, входной диапазон,
выходной размах и реакция на ёмкостную нагрузку могут различаться.
Высокая GBW сама по себе не гарантирует улучшения устройства.

\section{Питание: от вторичной обмотки до выводов NE5532}

\subsection{Что означают плюс, минус и GND}

В проекте три узла питания: +12 В, GND и −12 В. Все напряжения отсчитываются
от GND; между шинами +12 и −12 В получится 24 В.
Отрицательная шина не заменяет сигнальную землю. Обычный источник 12 В с двумя
контактами не создаёт одновременно эти три потенциала.

Двухполярное питание позволяет обрабатывать сигнал около нуля без искусственной
середины половины однополярного питания. При этом NE5532 всё равно имеет ограничения
входного синфазного диапазона и выходного размаха. Отступ от шин проверяют по
паспорту и измерению, а не по одному только знаку питания.

\subsection{Средняя точка и двухполупериодное выпрямление}

В \file{analog_power_lm317_lm337.tex} используются секции 5–6 и 5′–6′
трансформатора «Мелодия-103М». Соединение 6–6′ образует среднюю точку
только при требуемой фазировке. Для выбранного варианта сначала подтверждают
около 14–15 В AC на каждой секции и около 28–30 В между крайними выводами.
В \file{docs/POWER.md} напряжения предварительные; измерением они пока
не подтверждены. Переход на другую ревизию трансформатора требует пересчёта.

Один мост подключён к крайним выводам, а реальная средняя точка — к GND
конденсаторов. Оба плеча заряжаются в каждом полупериоде:
\begin{center}
\begin{tabularx}{\linewidth}{@{}Ycc@{}}
\toprule
Знаки крайних выводов относительно средней точки & Плюсовое плечо & Минусовое плечо\\
\midrule
5 положительный, 5′ отрицательный & VD1 & VD4\\
5 отрицательный, 5′ положительный & VD2 & VD3\\
\bottomrule
\end{tabularx}
\end{center}
Для сети 50 Гц основная частота пульсаций каждого плеча составляет 100 Гц.
Наличие средней точки не делает выпрямитель однополупериодным.

После накопительных конденсаторов уровень оценивается через амплитуду
полувторичной обмотки: $V_{\mathrm{raw}}\approx\sqrt2V_{\mathrm{AC,half}}-V_D$.
Это приближение для одного плеча. Падения на диодах, сопротивления обмоток,
изменение сети и зарядные импульсы требуют запаса. Суммарное напряжение
между положительной и отрицательной шиной примерно вдвое больше напряжения
одного плеча относительно GND.

\subsection{Накопительный конденсатор и CRC}

Первый конденсатор заряжается импульсами и питает нагрузку между ними.
Грубая оценка размаха пульсаций
\[
 \Delta V_{pp}\approx\frac{I}{f_{\mathrm{ripple}}C}.
\]
Три корпуса NE5532 в расчётном максимуме могут потребовать до 48 мА тока покоя
по указанному в даташите TI пределу 16 мА на корпус. Внешняя нагрузка источника
заложена 60 мА на плечо. Ещё около 11.4 мА берут задающие делители стабилизаторов.
Итого около 71.4 мА проходит через каждый CRC-резистор.

Для 1000 мкФ и 100 Гц оценка даёт примерно 0.71 В от пика до пика.
На 22 Ом падает около 1.57 В и выделяется примерно
\[
 P_R=I^2R\approx0.0714^2\cdot22=0.112\,\text{Вт}.
\]
В схеме стоят резисторы 1 Вт. Вместе со следующими 470 мкФ каждый образует
CRC-фильтр с характерным произведением $22\Ohm\cdot470\uF\approx10.3$ мс.
Номинал фильтра не заменяет проверку минимального напряжения на входах LM317/337.

Ток зарядки первого конденсатора существенно отличается от среднего тока
нагрузки. Его петля проходит через обмотку, диод и накопительный конденсатор.
Эту петлю делают компактной и отделяют от входных цепей фонокорректора.

\subsection{Что делают LM317D2T и LM337D2T}

Стабилизаторы уменьшают влияние колебаний входного напряжения, но им нужен
перепад между входом и выходом. В проекте желательно иметь минимум
около 15.5 В по модулю на входах при выходе ±12 В. Проверяется минимум
пульсирующего напряжения под нагрузкой, а не только среднее число мультиметра.

Задающий делитель использует около 1.25 В между OUT и ADJ.
При резисторе 110 Ом ток получается $1.25/110\approx11.4$ мА.
Он обеспечивает необходимую минимальную нагрузку и должен учитываться
в падении на фильтрах, нагреве и токе трансформатора.

При $|V_{\mathrm{in}}|=19$ В и токе 71.4 мА каждый регулятор рассеивает
примерно $(19-12)\cdot0.0714=0.50$ Вт. При 21 В — около 0.64 Вт.
Для корпуса D²PAK тепло отводится через площадку и медь платы.
Оценка по табличному тепловому сопротивлению не заменяет проверку реальной
платы внутри корпуса проигрывателя.

У выбранных мощных исполнений распиновка различается:
LM317 — ADJ 1, OUT 2, IN 3; LM337 — ADJ 1, IN 2, OUT 3.
Теплоотводящая площадка соединена с выводом 2. Поэтому полигон LM317 находится
на +12 В, а полигон LM337 — на нестабилизированном отрицательном входе.
Они не являются общим GND\cite{onsemi-317,onsemi-337}.

Для устойчивости LM337 важен выходной конденсатор и его свойства.
В нашей обвязке используется обычный алюминиевый электролит 10 мкФ;
произвольная замена его керамикой или полимером требует проверки по даташиту.
Защитные диоды задают пути разряда конденсаторов при изменении питания.
Их полярности обращены между положительным и отрицательным стабилизаторами.

\subsection{Локальная развязка каждого корпуса}

У каждой NE5532 стоят 100 нФ от вывода 8 к местному GND и 100 нФ от вывода 4
к местному GND; рядом — электролиты 47 мкФ. Малые конденсаторы нужны
непосредственно у выводов. Большой конденсатор у входа питания не заменяет
низкоиндуктивный локальный путь для быстрых токов.

На отрицательной шине плюс электролита обращён к земле, минус — к −12 В.
В паре межплечевых конденсаторов ситуация другая: их соединённые минусы
вообще не являются землёй. Совпадение знака «минус» на корпусах не означает,
что эти узлы нужно объединить.

\section{Разводка: карта токов важнее формы полигонов}

\subsection{Одна плата для стабилизаторов и предусилителя допустима}

На одной плате можно разместить LM317/337 и NE5532.
Для чувствительного входа следует выделить край платы рядом с разъёмами головки.
Элементы входов и обратной связи располагают компактно у U1/U2.
Стабилизаторы с теплоотводящей медью — в другой области; разностный U3 и
выходные разъёмы отделяют от входных проводов.

Мост и накопительные конденсаторы удобно оставить отдельным узлом.
Если они всё же находятся на общей плате, зарядные петли ограничивают её
силовой частью. Сам факт общего стеклотекстолита не создаёт неисправность;
важны электрическая и магнитная связь между конкретными контурами.

\subsection{GND на схеме идеален, на плате — импеданс}

Условные 100 мА через общий участок 10 мОм создают 1 мВ.
Для головки с расчётными 4 мВ это уже значимая величина.
На высокой частоте добавляется индуктивность, хорошо знакомая по ВЧ-платам.

Поэтому возвраты выпрямителя направляют к общей точке накопительных конденсаторов.
Возвраты нагрузки и делителей регуляторов прокладывают с учётом этих токов,
а входную сигнальную опору не используют как проходной участок силовой цепи.
Точка звезды — описание организации токов, а не требование стянуть
каждый ВЧ-развязывающий конденсатор длинным проводом в один угол корпуса.

Сплошная опорная медь в аналоговой области часто удобна.
Произвольное разрезание земли может увеличить площадь возвратных контуров.
Функциональное разделение областей платы и электрический разрыв GND — разные
решения. Теплоотводящие полигоны стабилизаторов дополнительно проверяются по
потенциалам, а не по внешнему сходству с полигоном земли.

\subsection{Экран, корпус, общий провод и защитное заземление}

Экран перехватывает электрическую помеху и соединяется по выбранной схеме
экранирования. Две жилы катушки остаются сигнальной парой.
У балансного входа экран не используется как замена L− или R−.
Соединение сигнальной земли с шасси организуют осознанно, учитывая остальные
подключённые приборы. USB к компьютеру способен добавить ещё один путь земли.

Защитное заземление сетевого оборудования не отключают для устранения фона.
Все описанные далее измерения относятся к низковольтной части устройства;
сетевая часть источника питания здесь не разрабатывается.

\subsection{Связь с механикой проигрывателя}

Головка является электромеханическим преобразователем. Геометрия установки,
прижим, состояние иглы, подшипники и работа двигателя влияют на результат
раньше электронного усиления. Их проверяют по механической конструкции и
паспорту головки. Электронная коррекция не компенсирует неправильный прижим
или повреждённую иглу.

Для электрических испытаний полезно отдельно сравнить состояния
«двигатель выключен» и «двигатель включён». Разница поможет обнаружить связь
через магнитное поле, вибрации или питание. Перемещение и поворот трансформатора
могут оказаться важнее небольшого изменения номинала конденсатора в тракте.

\section{Первое включение и измерения по шагам}

\subsection{Проверка монтажа и источника отдельно}

Перед подачей питания проверяют соответствие монтажу именно выбранного варианта:
балансный и небалансный имеют разные назначения деталей с похожими индексами.
Затем проверяют номера выводов, полярности конденсаторов, отсутствие замыкания
шин и состояние плавающей точки межплечевой пары.

Источник сначала проверяют без фонокорректора на контрольной нагрузке.
В схеме питания приведены резисторы 1 кОм / 0.25 Вт от каждой выходной шины
к GND для первоначальной настройки ±12 В. Для проверки полного бюджета
60 мА нужен другой эквивалент нагрузки с соответствующим рассеиванием;
контрольные 1 кОм потребляют лишь 12 мА на плечо.
Для 60 мА при 12 В подходит 200 Ом на каждое плечо относительно GND;
каждый резистор рассеивает 0.72 Вт, поэтому удобен номинал мощности 2 Вт.

При общем стендовом источнике полезно иметь ограничение тока и наблюдать
обе шины относительно GND. Сравнивают напряжения до и после CRC, входы
стабилизаторов и их выходы. Наличие ±12 В на холостом ходу не доказывает
достаточный запас под максимальной нагрузкой.

\subsection{Подать питание, пока головка и следующий усилитель отключены}

Для первоначального теста входы можно привести к местной земле через
предсказуемое малое сопротивление. Оставленный открытым вход способен собирать
наводки, и такой фон легко принять за собственный шум ОУ.
Проверяют ток потребления, напряжения питания у всех трёх корпусов и нагрев.

Далее измеряют постоянные напряжения на выходах U1/U2 и U3, затем на выходных
разъёмах после разделительных конденсаторов. На внутренних выходах некоторое
смещение допустимо. Сотни милливольт — повод сверить расчёт по входным токам
и монтаж, а выход, упершийся в шину, означает, что нормальный режим ещё не достигнут.
После выходного конденсатора и завершения переходного процесса уровень
относительно GND должен быть близок к нулю.

Затем осциллографом проверяют выходы и шины в широкой полосе.
Стабильный ноль на мультиметре не исключает ВЧ-генерацию.
Для измерения пульсаций питания используют короткое соединение земли щупа;
большая петля провода способна сама принять внешнюю помеху.

\subsection{Как подключать осциллограф}

У обычного заземлённого осциллографа земли каналов и корпуса BNC соединены.
Все земляные зажимы подключают к выбранному GND платы.
Их нельзя цеплять к L−, R−, отрицательной шине питания или плавающей середине
двух электролитов только потому, что это «минус» соответствующего сигнала.

Разность $X_+-X_-$ удобно наблюдать двумя щупами относительно GND и
математической операцией CH1−CH2. Входные ёмкости щупов и несовпадение каналов
ограничивают точность такого измерения. Для высокоомных входных узлов полезно
сначала оценить, насколько подключённый прибор меняет нагрузку.

\subsection{Генератор: ловушка нагрузки 50 Ом}

На генераторе надпись «1 В» может означать напряжение на согласованной нагрузке
50 Ом. Высокоомный вход получит другое напряжение, нередко вдвое больше.
Некоторые генераторы имеют отдельную настройку High-Z, влияющую на индикацию.
Поэтому действующее напряжение на входе платы нужно знать по измерению,
а не только по цифре на дисплее источника.

Для ослабления сигнала можно использовать делитель, например 1 кОм последовательно
и 10 Ом на землю. Отношение между \emph{измеренным напряжением перед делителем}
и его выходом близко к $10/(1000+10)$. Внутренние 50 Ом генератора влияют на
напряжение перед делителем; если же отсчитывать от внутреннего идеального
источника, они входят в знаменатель отдельно.

Калибровать удобнее на относительно большом уровне до делителя, проверив его
передаточную функцию и нагрузку. Длинные провода, измерительные приборы и
добавленные резисторы не должны превращать делитель в неизвестный фильтр.

\subsection{Проверить усиление на 1 кГц}

Для первого теста достаточно дифференциального входа около 0.5 мВ RMS.
Ожидаемый выход балансного варианта на 1 кГц — около 50.36 мВ RMS.
При 4 мВ RMS ожидается около 0.403 В RMS, но на больший уровень переходят
после проверки линейности и отсутствия генерации.

Для первичной проверки балансного входа без головки допустим односторонний
стендовый сигнал: генератор через ослабитель на L+, L− к опоре стенда через
согласованный малый импеданс. Между входами остаётся известное напряжение.
Такой тест проверяет передачу сигнала при наличии также небольшой синфазной
составляющей, но не доказывает CMRR и не задаёт окончательную проводку головки.
Для точной проверки режим возбуждения учитывают явно.

Правый канал проверяют тем же источником и тем же приспособлением.
Сравнение каналов при разных уровнях генератора теряет смысл.

\subsection{Снять АЧХ и не перегрузить низкие частоты}

При постоянном входе 0.5 мВ RMS для номинальной идеальной модели,
нагрузки 47 кОм и нулевого импеданса источника получаются:
\begin{center}
\IfFileExists{build/phono_guide/measurements.tex}{%
  \input{build/phono_guide/measurements.tex}%
}{\begin{tabular}{@{}rrrr@{}}\toprule
  Частота, Гц & Усиление, дБ & Усиление, В/В & Выход, мВ RMS\\\midrule
  20 & 59.27 & 919.22 & 459.61\\
  1000 & 40.06 & 100.72 & 50.36\\
  20000 & 20.44 & 10.52 & 5.26\\\bottomrule\end{tabular}}
\end{center}
Эта таблица — ориентир отладки, а не допуск готового прибора.
Для оценки формы АЧХ вычисляют
\[
 \varepsilon(f)=20\log_{10}\left|
 \frac{H_{\mathrm{meas}}(f)}{H_{\mathrm{meas}}(1\,\text{кГц})}
 \frac{H_{\mathrm{RIAA}}(1\,\text{кГц})}{H_{\mathrm{RIAA}}(f)}\right|.
\]
Измерение полного тракта «головка плюс корректор» — другая задача:
при нём источник уже не имеет нулевого импеданса.

Обратная RIAA-цепь позволяет получить почти плоский отклик пары
«испытательный фильтр плюс корректор». Но сама цепь тоже должна быть проверена.
Её точность, нагрузка, уровень и ограниченная рабочая полоса входят в результат.
При цифровом формировании испытательного синуса можно вместо отдельного фильтра
задавать амплитуду по частоте, а затем всё равно калибровать аналоговый выход.

\subsection{Почему прямоугольник на выходе не обязан оставаться прямоугольником}

Прямоугольный сигнал содержит набор гармоник, а фонокорректор усиливает их
по-разному. Поэтому форма изменится даже у идеально линейной схемы.
Привычный критерий импульсного тракта «повторить фронт без изменения»
не является целью RIAA-каскада.

Проверка переходного процесса всё равно полезна для поиска возбуждения,
но интерпретировать её нужно через ожидаемую передаточную функцию.
АЧХ измеряют синусом или другим калиброванным методом; устойчивость дополнительно
проверяют в широкой полосе. Высокочастотный паразитный звон не следует
автоматически объяснять полезной коррекцией.

\subsection{Проверить общий режим, шум и перегрузку}

Для отдельной синфазной проверки головку отключают и подают одинаковый сигнал
на оба сигнальных входа одного канала. Начальный уровень порядка 0.1 В RMS
остаётся удобным малосигнальным испытанием; фактические внутренние уровни всё
равно контролируют. В идеале выход равен нулю, на макете он ограничен ошибками
сети, источника и приборов.

Если $A_d=V_{o,d}/V_{i,d}$ и $A_{cm}=V_{o,cm}/V_{i,cm}$, то
$CMRR=20\log_{10}|A_d/A_{cm}|$ при одинаковой частоте.
Разница амплитуд дифференциального и синфазного тестов должна учитываться.
Результат ниже шумового пола прибора даёт лишь нижнюю границу CMRR.

Шум сначала сравнивают с известным источником или эквивалентом головки.
Затем подключают реальную головку и отдельно включают двигатель.
Фон, появившийся только после подключения кабеля или USB, имеет полезную
диагностическую ценность. Для проверки перегрузки повышают синусоидальный
уровень и наблюдают все внутренние выходы, а не только выходной разъём.

Искажения измеряются вместе с остаточными искажениями генератора и анализатора.
Если их собственный предел выше ожидаемого результата предусилителя,
уверенное численное заключение о THD не получится.

\section{Как устроены расчётные скрипты проекта}

\subsection{Файлы, единицы и запуск}

\begin{tabularx}{\linewidth}{@{}L{0.49\linewidth}Y@{}}
\toprule
Файл & Назначение\\
\midrule
\file{config.toml} & Номиналы и условия расчёта. R в омах, C в фарадах,
время в секундах, частота в герцах.\\
\file{scripts/calculate_balanced_riaa.py} & Балансный вариант, три ОУ на канал.\\
\file{scripts/calculate_unbalanced_riaa.py} & Исходный вариант, один ОУ на канал.\\
\file{scripts/riaa_common.py} & Общие узловые уравнения, проверка,
таблицы и графики.\\
\file{pyproject.toml} & Зависимости Python; установка в существующую .venv.\\
\file{scripts/render_phono_guide.py} & Иллюстрации и числовая таблица этого
руководства. Читает модели и пишет только новые результаты руководства.\\
\bottomrule
\end{tabularx}

Из корня проекта:
\begin{verbatim}
.venv/bin/python -m pip install -e .
.venv/bin/python scripts/calculate_balanced_riaa.py
.venv/bin/python scripts/calculate_unbalanced_riaa.py
\end{verbatim}
У каждого основного скрипта есть \file{--config PATH} и \file{--output-dir DIR}.
По умолчанию CSV, JSON, PNG и PDF сохраняются в \file{build/balanced_riaa/}.
В JSON записывается и конфигурация, с которой получен результат.

Для повторной сборки этого руководства:
\begin{verbatim}
.venv/bin/python scripts/render_phono_guide.py
latexmk -xelatex -outdir=build/phono_guide docs/PHONO_GUIDE.tex
\end{verbatim}
После сборки итоговый \file{build/phono_guide/PHONO_GUIDE.pdf} можно
скопировать в \file{docs/PHONO_GUIDE.pdf}. Без повторного запуска расчёта
старые рисунки сохраняют предыдущий набор параметров.

\subsection{Формула и независимый узловой расчёт}

Первый способ вычисляет передаточную функцию из аналитического выражения.
Второй строит комплексную матрицу проводимостей по законам Кирхгофа.
Строки выходов ОУ задают соотношение $V_o=A(s)(V_+-V_-)$.
Идеальная модель получается при $1/A=0$.

Для балансного варианта два конденсатора по 1000 мкФ представлены раздельно,
а их общая точка — самостоятельный плавающий узел.
После решения проверяется совпадение с формулой. Дополнительно на оба входа
подаётся одинаковое напряжение и проверяется идеальное подавление общего режима.

Разница на уровне численной погрешности подтверждает согласованность двух
способов расчёта одной модели. Она не подтверждает, что в модель внесены
все свойства реальных компонентов. Например, одинаковый ошибочный номинал
в конфигурации будет одинаково использован обоими расчётами.

\subsection{Какие вопросы модели не решают}

В текущих расчётах нет импеданса катушки, кабеля, собственных шумовых источников,
смещения, входных токов, утечек и нелинейности электролитов. Нет ограничений
выходного напряжения, slew rate, детального выходного импеданса ОУ и паразитных
полюсов платы. Линейная модель не создаёт гармоник и сама по себе не считает THD.

Кривая с однополюсным ОУ показывает влияние конечного усиления и GBW на
линейную передачу. Её нельзя считать SPICE-макромоделью NE5532 или доказательством
устойчивости. Идеальная кривая показывает, насколько правильно подобраны
номиналы активного фильтра при принятых допущениях.

Каждая кривая ошибки нормирована по своему усилению на 1 кГц.
Поэтому совпадение ошибок в нуле на этой частоте не означает равенства
абсолютного усиления. Для сравнения одновременно смотрят таблицу коэффициентов.

Изменение \file{config.toml} меняет расчёт, но не подписи в TeX.
При пересчёте реального устройства обновляют также схему, перечень компонентов
и протокол измерений. Само это руководство является отдельным пояснением,
а не автоматической системой изменения проекта.

\section{Что считать хорошим результатом первого макета}

Для начала полезнее получить воспроизводимый рабочий тракт, чем назначать ему
неподтверждённый паспорт с очень малой цифрой THD. Результат оценивается
по нескольким независимым признакам:
\begin{enumerate}
\item Питание имеет правильные полярности, держит нагрузку и не выводит
стабилизаторы за допустимые режимы.
\item ОУ находятся в линейном режиме, не генерируют и не перегреваются.
\item Усиление и форма АЧХ соответствуют расчёту в пределах объяснимых
допусков, а левый и правый каналы согласованы.
\item Известны условия измерения шума, вклад двигателя, трансформатора,
кабеля и подключённой аппаратуры.
\item Есть измеренный запас до перегрузки на нескольких частотах,
включая контроль промежуточных каскадов.
\item Реальная головка и механика проигрывателя корректно установлены,
после чего прослушивание не противоречит измеренным результатам.
\end{enumerate}

При сравнении вариантов следует выровнять уровень на 1 кГц и сохранить
одинаковые нагрузку, источник и положение кабелей. Более громкий вариант
легко воспринимается как более выразительный. Сравнение с включённым
двигателем и реальным кабелем проверяет другие свойства, чем сравнение
двух плат с закороченными входами на лабораторном столе.

Внешний вид скрытой платы действительно вторичен. При этом аккуратный монтаж
имеет электрический смысл: короткие обратные связи, понятные точки измерения,
разнесённые зарядные петли, надёжные разъёмы и возможность повторить результат.
Большие электролиты сами по себе не ухудшают схему только из-за размера.
Удачное размещение их выводов и возвратных токов полезнее декоративной симметрии.

\section{Если результат не совпал с ожиданием}

\begin{longtable}{@{}L{0.23\linewidth}L{0.35\linewidth}L{0.34\linewidth}@{}}
\toprule
Симптом & Возможная причина & Что проверить первым\\
\midrule\endhead
Выход у шины питания & Нет обратной связи, перепутаны выводы, вход
не имеет DC-пути, ошибочная полярность питания & Напряжения на выводах 8/4,
потенциалы обоих входов и непрерывность ООС.\\
Мультиметр показывает норму, ОУ греется & ВЧ-генерация или лишняя нагрузка &
Осциллограф в широкой полосе, локальная развязка, выходные кабели.\\
Уровень вдвое отличается & Настройка генератора 50 Ом/High-Z или перепутаны
RMS, пик и пик-пик & Фактическую амплитуду до делителя и определение входного сигнала.\\
Сильно ослаблены басы & Малая эффективная ёмкость, две ёмкости в серии
не учтены, другое сопротивление нагрузки & Номиналы и эквивалентные RC-произведения,
особенно 1000+1000 мкФ = 500 мкФ.\\
Ослаблены высокие частоты & Добавлен второй полюс 75 мкс, ошибочен C,
вход нагружен лишней ёмкостью & Наличие только предусмотренной ВЧ-коррекции,
номиналы нФ/мкФ и нагрузку головки.\\
Два канала заметно различаются & Ошибка детали, допуски RIAA, разная нагрузка &
Один и тот же источник, перечень деталей и парные измерения R/C.\\
Фон появляется с головкой & Кабель принимает помеху, нежелательная связь
L−/R− с шасси, источник стал индуктивным & Проводку, экраны, контакты головки
и сравнение с корректным эквивалентом.\\
Фон усиливается с двигателем & Магнитная, механическая или проводная связь &
Расположение узлов, трансформатор, питание двигателя, состояние механики.\\
Фон появляется с USB & Новый путь общего провода через компьютер &
Все соединения сигнальной земли и заземления приборов.\\
После U3 уровень мал, но слышны искажения & Перегрузка U1/U2 или сам источник
искажён & Оба внутренних выхода и сигнал генератора под реальной нагрузкой.\\
Щелчок при подключении или включении & Заряд разделительных ёмкостей,
переход смещений, неправильный выходной DC-режим & Постоянные уровни,
разрядные резисторы и установившийся режим до подключения нагрузки.\\
\bottomrule
\end{longtable}

Это список направлений поиска. Один симптом может иметь несколько причин.
Сначала меняют один фактор и фиксируют результат, затем переходят к следующему.
Замена нескольких деталей и проводов одновременно затрудняет вывод о механизме.

\section{Памятка по конкретному выбранному варианту}

\begin{itemize}
\item Основной балансный вариант использует три корпуса NE5532 на стерео.
L− и R− являются сигнальными проводами; экран и земля тонарма отдельны.
\item В каждом входном плече стоит 23.5 кОм к GND; между концами катушки
получается 47 кОм в рабочей полосе.
\item Между инвертирующими входами каждого канала стоят 182 Ом и два
электролита по 1000 мкФ / 25 В, минусами вместе. Средняя точка плавающая.
\item Частотно-зависимая ООС входной пары использует 90.9 кОм,
10 кОм и 31.5 нФ. Её изменение требует пересчёта.
\item $X_{L\pm}$ и $X_{R\pm}$ ещё не являются готовыми линейными выходами:
в этих точках нет спада 75 мкс.
\item В U3 четыре согласованных 10 кОм и два согласованных 7.5 нФ
на каждый канал выполняют вычитание и окончательную ВЧ-коррекцию.
\item Прежний выходной фильтр 750 Ом / 100 нФ к этому варианту не добавляется.
После U3 остаются 100 Ом, 4.7 мкФ и 100 кОм.
\item Питание — ±12 В; бюджет источника — 60 мА внешней нагрузки на плечо
плюс примерно 11.4 мА его задающего делителя.
\item Площадка LM317D2T — +12 В, площадка LM337D2T — отрицательный
нестабилизированный вход. Это разные полигоны, не GND.
\item Расчётное усиление около 100.72 В/В на 1 кГц относится к напряжению
между входными сигнальными проводами. Оно не является коэффициентом
каждой отдельной секции и не постоянно по звуковой полосе.
\end{itemize}

\section{Короткий словарь аудиотерминов}

\begin{description}[style=nextline]
\item[Фонокорректор]
Предусилитель с частотной коррекцией для магнитной головки проигрывателя.
\item[PHONO и LINE]
PHONO — слабый сигнал головки, требующий коррекции. LINE — сигнал линейного
уровня для следующего тракта. Надписи на конкретном проигрывателе проверяют
по его схеме: встроенный корректор может уже выполнять преобразование.
\item[MM и MC]
Moving Magnet и Moving Coil — разные принципы головок с различными
уровнями и импедансами. Нагрузку и усиление одной конструкции нельзя
автоматически переносить на другую.
\item[RIAA]
Заданная форма частотной коррекции записи/воспроизведения пластинок.
\item[Балансный вход]
Вход, принимающий разность двух сигнальных проводов при согласованных
условиях их связи с общей опорой. Две жилы сами по себе ещё не задают
подавление синфазной помехи.
\item[CMRR и PSRR]
CMRR характеризует подавление синфазного сигнала на входах. PSRR — влияние
изменения питания на ошибку усилителя. Это разные пути возникновения помехи.
\item[Noise gain]
Передача эквивалентного входного шумового напряжения ОУ на выход.
Она не обязана совпадать с передачей внешнего полезного сигнала.
\item[Смещение и входной ток]
Постоянные ошибки реального ОУ. Отличаются от спектральной плотности
шума и могут ограничивать доступный размах даже без музыки.
\item[BP, NP и NP0]
BP/NP у электролита означают неполярное исполнение. NP0 у керамики —
обозначение температурного поведения, связанное с классом C0G.
Это не три названия одного типа большого конденсатора.
\item[Рокот и subsonic]
Рокот — низкочастотная составляющая, связанная с механикой и приводом.
Subsonic-фильтр ограничивает инфранизкочастотную часть. Его АЧХ должна
учитываться отдельно от чистой трёхвременной RIAA.
\item[A-взвешивание]
Частотное взвешивание при оценке шума. Оно меняет итоговую цифру и должно
быть явно указано в протоколе.
\item[DC-серво]
Медленная дополнительная петля для контроля постоянного уровня.
Может позволить иначе организовать разделение DC/AC, но требует отдельного
проектирования и не является уже установленным узлом нашей схемы.
\end{description}

\section{Граница с будущей оцифровкой}

Линейный выход фонокорректора — удобная точка для подключения следующего
аналогового или цифрового тракта. Для АЦП нужно заново проверить допустимый
размах, постоянное смещение, входную нагрузку и тип входа: дифференциальный
или относительно общей точки. Наличие слова «24 bit» не заменяет этих условий.

RIAA не является достаточным антиалиасинговым фильтром для произвольного
АЦП. Выбранный преобразователь может иметь собственный аналоговый входной
каскад, передискретизацию и цифровые фильтры, но его требования нужно
проверять отдельно. Сигнал 0.4 В RMS с отрицательной полуволной нельзя автоматически подключать
к однополярному входу микроконтроллера.

Питание цифровой части, передача данных по USB и общий провод с компьютером
способны изменить помеховую обстановку. Поэтому полезно сначала измерить
самостоятельный аналоговый тракт, а затем повторить измерения после добавления
цифровой платы. Такой порядок отделяет свойства фонокорректора от свойств
нового интерфейса.

\clearpage
\phantomsection
\addcontentsline{toc}{section}{Справочные документы}
\begin{thebibliography}{99}
\bibitem{ti-ne5532}
Texas Instruments. \emph{NE5532x, SA5532x Dual Low-Noise Operational Amplifiers},
SLOS075K. Локально: \file{docs/datasheet/OpAmp/ne5532.pdf}.
\href{https://www.ti.com/lit/ds/symlink/ne5532.pdf}{Официальный даташит}.
В частности, таблицы распиновки, токов и раздел Layout.

\bibitem{onsemi-ne5532}
onsemi. \emph{NE5532/D}. Локально:
\file{docs/datasheet/OpAmp/NE5532-D.PDF}.
При выборе конкретного производителя используют именно его ограничения,
а не удобную смесь типовых параметров разных версий.

\bibitem{onsemi-audio}
onsemi. \emph{Audio Circuits Using the NE5532/4}, AND8177/D.
Локально: \file{docs/datasheet/OpAmp/AND8177-D.PDF}.
\href{https://www.onsemi.com/download/application-notes/pdf/and8177-d.pdf}{Официальная публикация}.
Рисунок 3 находится на печатной странице 2, третьей странице PDF.
В подписи верхней шины этого рисунка есть опечатка «−15 V»:
для NE5532 вывод 8 является положительным питанием.

\bibitem{adi-ina}
Analog Devices. \emph{MT-063: Basic Three Op Amp In-Amp Configuration}.
\href{https://www.analog.com/media/en/training-seminars/tutorials/MT-063.pdf}{Архитектура входной пары и разностного каскада}.

\bibitem{adi-bias}
Analog Devices. \emph{AN-937: Designing Amplifier Circuits: How to Avoid Common Problems}.
\href{https://www.analog.com/en/resources/app-notes/an-937.html}{Пути входных токов и связанные ошибки}.

\bibitem{adi-noise}
Analog Devices. \emph{MT-047: Op Amp Noise}.
\href{https://www.analog.com/media/en/training-seminars/tutorials/MT-047.pdf}{Модели источников шума и шум сопротивлений}.

\bibitem{cde}
Cornell Dubilier. \emph{Application Guide, Aluminum Electrolytic Capacitors},
раздел Non-Polar / Anti-series connection.
Локально: \file{docs/datasheet/capacitor/c04-appguide.pdf}.
\href{https://knowledge.electrochem.org/encycl/c04-appguide.pdf}{Копия руководства производителя}.

\bibitem{ti-caps}
Texas Instruments. \emph{Selecting Capacitors to Minimize Distortion in Audio Applications},
SLYT796A. \href{https://www.ti.com/lit/an/slyt796a/slyt796a.pdf}{Сравнение конденсаторов в конкретной сигнальной цепи}.

\bibitem{murata}
Murata. Ceramic Capacitors FAQ, acoustic noise / piezoelectric effect.
\href{https://www.murata.com/support/faqs/capacitor/ceramiccapacitor/char/0020}{Электромеханические свойства ферроэлектрической керамики}.

\bibitem{onsemi-317}
onsemi. \emph{LM317/D}.
\href{https://www.onsemi.com/pdf/datasheet/lm317-d.pdf}{Распиновка D²PAK и электрические ограничения}.

\bibitem{onsemi-337}
onsemi. \emph{LM337/D}. Локально:
\file{docs/datasheet/ldo/LM337-D.PDF}.
\href{https://www.onsemi.com/pdf/datasheet/lm337-d.pdf}{Распиновка и требования к выходному конденсатору}.
\end{thebibliography}

\end{document}
